\notechapter{N-20220501}{Lagrange Multipliers, Hessians, Bordered Hessians, and Degenerate Extrema}

\section{The full problem of constrained extrema}

Constrained optimization is not merely “Lagrange multipliers plus Hessian.”  A complete treatment must explain the multiplier geometrically, derive the stationarity condition, handle bordered Hessians and singular constraints, examine higher-order degeneracy, and connect the theory with numerical optimization.

\section{Lagrange multipliers: the basic idea}

Suppose we want to find extrema of
\[
  f(x_1,\ldots,x_n)
\]
under the constraint
\[
  g(x_1,\ldots,x_n)=0.
\]
The Lagrange function is
\[
  L(x_1,\ldots,x_n,\lambda)
  =f(x_1,\ldots,x_n)-\lambda g(x_1,\ldots,x_n).
\]
The stationarity equations are
\[
  \frac{\partial L}{\partial x_i}=0,
  \qquad i=1,\ldots,n,
  \qquad
  \frac{\partial L}{\partial \lambda}=0.
\]
Equivalently,
\[
  \nabla f=\lambda \nabla g,
  \qquad
  g=0.
\]
Geometrically, at a constrained extremum, the level surface of $f$ is tangent to the constraint surface.  Therefore their normal directions are parallel.  This explains why gradients are proportional.

The method can seem surprising because it turns a constrained problem into a stationarity system, and confusing because the multiplier appears to come from nowhere.  Geometrically, $\lambda$ is simply the scalar comparing the two normal vectors.

\section{Example: minimizing distance to a line}

Consider
\[
  f(x,y)=x^2+y^2
\]
under the constraint
\[
  x+y=1.
\]
Set
\[
  L(x,y,\lambda)=x^2+y^2-\lambda(x+y-1).
\]
The equations are
\[
  2x-\lambda=0,
  \qquad
  2y-\lambda=0,
  \qquad
  x+y=1.
\]
The first two give $x=y$, and therefore $x=y=1/2$.  This is the closest point of the line $x+y=1$ to the origin.

The example is simple but instructive.  For nonlinear constraints, the stationarity system can become hard to solve, and numerical methods such as Newton iteration may be necessary.  Lagrange multipliers produce equations, not automatic closed forms.

\section{Why the method is valid}

A more rigorous proof uses the implicit function theorem.  If $\nabla g(p)\ne0$, then near $p$ the constraint $g=0$ is a smooth hypersurface.  A tangent vector $v$ to the constraint satisfies
\[
  \nabla g(p)\cdot v=0.
\]
If $p$ is a constrained extremum of $f$, then the derivative of $f$ along every tangent direction must vanish:
\[
  \nabla f(p)\cdot v=0
  \qquad \text{for all tangent }v.
\]
Thus $\nabla f(p)$ is perpendicular to the tangent space.  The normal space is spanned by $\nabla g(p)$, so
\[
  \nabla f(p)=\lambda \nabla g(p).
\]
For several constraints $g_1=\cdots=g_m=0$, the condition becomes
\[
  \nabla f=\sum_{i=1}^m \lambda_i \nabla g_i,
\]
assuming the constraint gradients are independent.

This is why singular constraints require care: if the constraint gradients are linearly dependent, the usual bordered-Hessian test may fail.

\section{Hessian matrix and the second derivative test}

For a twice differentiable function $f(x_1,\ldots,x_n)$, the Hessian matrix at $x_0$ is
\[
  H_f(x_0)=\left(\frac{\partial^2 f}{\partial x_i\partial x_j}(x_0)\right)_{i,j}.
\]
If the mixed partials are continuous, the Hessian is symmetric.  This lets us use the linear algebra of symmetric matrices.

At an unconstrained critical point $\nabla f(x_0)=0$:
\begin{enumerate}[(i)]
  \item if $H_f(x_0)$ is positive definite, $x_0$ is a local minimum;
  \item if $H_f(x_0)$ is negative definite, $x_0$ is a local maximum;
  \item if $H_f(x_0)$ is indefinite, $x_0$ is a saddle point;
  \item if $H_f(x_0)$ is singular or semidefinite, the second derivative test is inconclusive.
\end{enumerate}

The proof comes from Taylor expansion:
\[
  f(x_0+h)=f(x_0)+\frac12 h^T H_f(x_0)h+o(\|h\|^2)
\]
when the gradient term vanishes.  Thus the sign of the quadratic form controls the local shape.

\section{Sylvester criterion and eigenvalues}

For a real symmetric matrix, positive definiteness can be tested by eigenvalues or by leading principal minors.  Sylvester's criterion says that $A$ is positive definite iff all leading principal minors are positive.  Negative definiteness has alternating signs.  Eigenvalues give an equivalent test: positive definite means all eigenvalues are positive; indefinite means there are both positive and negative eigenvalues.

This is a direct linear-algebra connection.  A Hessian is not just an array of second derivatives; it is a quadratic form.  Its eigenvectors are principal directions of curvature, and its eigenvalues measure curvature along those directions.

\section{Bordered Hessian}

For constrained extrema, the ordinary Hessian of $f$ alone is not enough because admissible motion is restricted to directions tangent to the constraint.  The bordered Hessian incorporates constraint-gradient information into the second derivative test.

For constraints $g_i(x)=0$, define the Lagrangian
\[
  L=f-\sum_{i=1}^m \lambda_i g_i.
\]
A common bordered Hessian has the block form
\[
  B=
  \begin{pmatrix}
  0_{m\times m} & (Dg)^T\\
  Dg & H_x L
  \end{pmatrix},
\]
where $Dg$ records gradients of the constraints and $H_xL$ is the Hessian with respect to the original variables.  Different texts use slightly different sign conventions and ordering of rows, so one must check the convention before applying a sign rule.

The intuitive meaning is: the border records which directions are forbidden by the constraints, and the Hessian block records curvature of the Lagrangian in the remaining directions.

\section{Example revisited with a bordered Hessian}

For the simple constraint $x+y=1$, the constraint gradient is
\[
  \nabla g=(1,1).
\]
If $f(x,y)=x^2+y^2$, then
\[
  H_f=
  \begin{pmatrix}
  2&0\\
  0&2
  \end{pmatrix}.
\]
The tangent direction to $x+y=1$ is spanned by $(1,-1)$.  Along that direction,
\[
  (1,-1)H_f(1,-1)^T=4>0,
\]
so the constrained critical point is a constrained minimum.

This tangent-space calculation is often clearer than memorizing signs of bordered determinants.  The bordered Hessian is a compact determinant method for doing the same restriction systematically.

\section{Degenerate cases and higher-order tests}

When the Hessian is singular or the bordered-Hessian test is inconclusive, second-order information is not enough.  One must examine higher-order Taylor terms:
\[
  f(x_0+h)=f(x_0)+P_k(h)+o(\|h\|^k),
\]
where $P_k$ is the first nonzero homogeneous term after the constant term.  If $P_k$ has a fixed sign, one may still obtain an extremum; if it changes sign, one gets a saddle-type point.

A standard example is $f(x,y)=x^4+y^4$, whose Hessian at the origin is zero, but the origin is still a strict local minimum.  Conversely, $f(x,y)=x^4-y^4$ also has zero Hessian at the origin but has a saddle.  Thus degeneracy forces us to look beyond quadratic approximation.

\section{Saddle points and geometric visualization}

A saddle point is a critical point that is neither a local maximum nor a local minimum.  If the Hessian has both positive and negative eigenvalues, then some directions curve upward and others downward.  Contour plots are useful for detecting saddle behavior because level curves reveal local shape far more clearly than raw formulas.

The classical picture is $f(x,y)=x^2-y^2$.  Along the $x$-axis it increases; along the $y$-axis it decreases.  The origin is therefore a saddle.

\section{Optimization beyond equality constraints}

Kuhn--Tucker conditions, penalty methods, and interior-point methods belong to constrained optimization with inequalities.  If constraints include
\[
  h_j(x)\le0,
\]
then Lagrange multipliers are supplemented by nonnegativity and complementary slackness:
\[
  \mu_j\ge0,
  \qquad
  \mu_jh_j(x)=0.
\]
This says that an inequality constraint contributes a force only when it is active.  This is the natural extension of the Lagrange multiplier idea.

Penalty methods instead replace constraints by large penalty terms, for example
\[
  f(x)+\rho \sum_j h_j(x)^2.
\]
Interior-point methods keep the iterate inside the feasible region using barrier functions.  All of these methods share the same philosophy: convert a constrained problem into a system that can be treated with unconstrained or nearly unconstrained tools.

\section{Global topology and extrema}

Global topology enters naturally.  Local Hessian information controls a neighborhood of a critical point, while global existence of maxima and minima depends on compactness.  A continuous function on a compact set attains maximum and minimum.  Thus topology enters before calculus even begins.

There are also deeper links: Morse theory studies how critical points of smooth functions reflect the topology of a manifold.  The Hessian at a nondegenerate critical point gives the Morse index, and the collection of critical points helps reconstruct the space.  This is the high-level form of the principle that extrema are connected with global structure.

\section{Computation and MATLAB}

For high-dimensional functions, symbolic or numerical software can compute gradients, Hessians, eigenvalues, and bordered Hessians.  A typical workflow is:
\begin{enumerate}[(i)]
  \item define $f$ and constraints $g_i$;
  \item solve $\nabla f=\sum \lambda_i\nabla g_i$ together with $g_i=0$;
  \item compute the Hessian of the Lagrangian;
  \item restrict to the tangent space or use a bordered Hessian;
  \item inspect eigenvalues or determinant signs.
\end{enumerate}
This computational workflow is mathematically meaningful: software does not replace theory, but automates the linear algebra after the theoretical conditions have been derived.


\section{Second variation on the tangent space}

For constrained extrema, the cleanest conceptual test is to restrict the Hessian of the Lagrangian to tangent directions.  If the constraint is \(g(x)=0\), then tangent vectors \(v\) satisfy
\[
  \nabla g(x_0)\cdot v=0.
\]
The second variation is
\[
  v^T H_xL(x_0,\lambda)v.
\]
If this quadratic form is positive definite on the tangent space, the point is a constrained local minimum; if negative definite, a constrained local maximum.

This explains what the bordered Hessian is doing: it is a determinant device for testing the same restricted quadratic form.

\section{From constrained extrema to second variation and Morse data}

The transition from calculus to optimization and differential geometry is structural.  Lagrange multipliers are normal-vector equations; Hessians are quadratic forms; bordered Hessians restrict second variation to tangent directions; degenerate cases require higher-order terms; inequality constraints lead to KKT theory; compactness guarantees global extrema; and Morse theory relates critical points to topology.

\section{Lagrange multipliers as cotangent annihilators}

For constraints \(g_1=\cdots=g_k=0\), the tangent space consists of directions annihilated by all \(dg_i\).  At a constrained extremum,
\[
  df\in \operatorname{span}(dg_1,\ldots,dg_k).
\]
This is the coordinate-free form of Lagrange multipliers.

\section{Bordered Hessian and restricted second variation}

The second-order test should be applied only to tangent directions of the constraint.  The bordered Hessian is a coordinate device for computing the quadratic form of the Lagrangian on that tangent space.

If the restricted quadratic form is positive definite, the point is a constrained local minimum; if negative definite, a constrained local maximum.

\section{Degenerate extrema and higher order}

When the second variation has zero directions, the Hessian test is inconclusive.  One must examine higher-order terms or use normal forms.  This is the optimization analogue of a multiple root in algebra: first-order and second-order information may not separate the cases.

\section{First and second variations under constraints}

Constrained extrema are governed by tangent and cotangent spaces.  Multipliers express first-order normality, bordered Hessians test second-order behavior along constraints, and degeneracy signals the need for higher-order local geometry.
